% !TeX root = ../main_theorems_multifile.tex
% Paper 1 appendix: algebraic gauge structure (T1). Locus geometry lives in Paper 2.

\section{Gauge structure in Karlsson's family}
\label{sec:proof-gauge}

Write Karlsson's three-parameter family as $K_6(\theta,\phi,\lambda)$ with block $A(\theta,\phi)$ and $z_1=e^{i\lambda}$ determining the remaining M\"obius blocks (as in Section~\ref{sec:background}).
Throughout, $H(\theta,\phi,\lambda):=K_6(\theta,\phi,\lambda)$.

\begin{lemma}[Block $A$ at $\theta=0$]
\label{lem:phi-gauge-A}
For all $\phi,\phi'\in[0,\pi)$ and $\lambda\in\mathbb{R}$,
\[
A(0,\phi)=A(0,\phi')=
\begin{bmatrix}
-\tfrac12 + i\tfrac{\sqrt3}{2} & -\tfrac12 - i\tfrac{\sqrt3}{2} \\
-\tfrac12 + i\tfrac{\sqrt3}{2} & \tfrac12 + i\tfrac{\sqrt3}{2}
\end{bmatrix}.
\]
\end{lemma}

\begin{proof}
Substitute $\sin\theta=0$ in the Karlsson formulas
\[
A_{11}=-\tfrac12 + i\tfrac{\sqrt3}{2}(\cos\theta+e^{-i\phi}\sin\theta),\quad
A_{12}=-\tfrac12 + i\tfrac{\sqrt3}{2}(-\cos\theta+e^{i\phi}\sin\theta).
\]
All $\phi$-dependent terms vanish, giving the displayed matrix.
Unitarity $AA^\dagger=2I$ follows by direct multiplication.
\end{proof}

\begin{lemma}[$\phi$-gauge for $H$ at $\theta=0$]
\label{lem:L1}
For all $\phi,\phi',\lambda$: $H(0,\phi,\lambda)=H(0,\phi',\lambda)$ entrywise.
\end{lemma}

\begin{proof}
At $\theta=0$, Lemma~\ref{lem:phi-gauge-A} gives $A$ and $B=-F_2-A$ independent of $\phi$.
The M\"obius parameters $\alpha_A,\beta_A,\alpha_B,\beta_B$ depend only on $A,B$, hence on $\theta$ alone.
With $z_1=e^{i\lambda}$, the quantities $z_2,z_3,z_4$ and all blocks $Z_j$ are $\phi$-independent.
Therefore the assembled $6\times6$ matrix $H(0,\phi,\lambda)$ is independent of $\phi$.
\end{proof}

\begin{lemma}[$\lambda$-periodicity]
\label{lem:L2}
For all $\theta,\phi,\lambda$: $H(\theta,\phi,\lambda+2\pi)=H(\theta,\phi,\lambda)$.
\end{lemma}

\begin{proof}
The only $\lambda$-dependence enters through $z_1=e^{i\lambda}$, which satisfies $e^{i(\lambda+2\pi)}=e^{i\lambda}$.
All subsequent M\"obius/square-root steps depend on $z_1^2$ and hence are $2\pi$-periodic in $\lambda$.
\end{proof}

\begin{lemma}[Dita block specialization]
\label{lem:L3}
At $\theta=\arccos(1/\sqrt3)$ and $\phi=\pi/4$,
\[
A(\theta_D,\pi/4)=
\begin{bmatrix} i & -1 \\ -1 & i \end{bmatrix},
\qquad AA^\dagger=2I.
\]
\end{lemma}

\begin{proof}
With $\cos\theta=1/\sqrt3$, $\sin\theta=\sqrt{2/3}$, and $\phi=\pi/4$,
\[
A_{11}=-\tfrac12+i\tfrac{\sqrt3}{2}\Bigl(\tfrac{1}{\sqrt3}+\tfrac{1-i}{\sqrt3}\Bigr)=i,\quad
A_{12}=-\tfrac12+i\tfrac{\sqrt3}{2}\Bigl(-\tfrac{1}{\sqrt3}+\tfrac{1+i}{\sqrt3}\Bigr)=-1.
\]
Hermitian structure $A=\begin{pmatrix}A_{11}&A_{12}\\ \overline{A_{12}}&-\overline{A_{11}}\end{pmatrix}$ yields the matrix shown.
Direct computation gives $AA^\dagger=2I$.
\end{proof}

\begin{lemma}[Matrix inequality in $\lambda$ at Dita]
\label{lem:L4-matrix}
Let $\theta_D=\arccos(1/\sqrt3)$, $\phi=\pi/4$, and write $H(\lambda):=H(\theta_D,\phi,\lambda)$.
If $\lambda\not\equiv\lambda'\pmod{2\pi}$, then $H(\lambda)\neq H(\lambda')$ as matrices.
\end{lemma}

\begin{proof}
In the block assembly, the one-based entry $H(\lambda)_{2,3}$ is $(Z_1)_{2,1}=z_1=e^{i\lambda}$. Therefore equality of $H(\lambda)$ and $H(\lambda')$ implies $e^{i\lambda}=e^{i\lambda'}$, hence $\lambda\equiv\lambda'\pmod{2\pi}$.
\end{proof}

\begin{lemma}[CHM class structure on the Dita $\lambda$-circle]
\label{lem:L4}
\label{lem:dita-chm-pi}
Let $\theta_D=\arccos(1/\sqrt3)$, $\phi=\pi/4$, and $H(\lambda):=H(\theta_D,\phi,\lambda)$.
\begin{enumerate}
    \item \emph{$\pi$-periodicity of CHM class.} Dephased alignment over $720$ column permutations and diagonal phases yields residuals $\sim10^{-15}$ for the pairs $(\lambda,\lambda+\pi)$ with $\lambda\in\{0,\pi/3,2\pi/3\}$ (\repo{chm\_equivalence\_residual}).
    Thus $H(\lambda)$ and $H(\lambda+\pi)$ are CHM-equivalent on this slice.
    \item \emph{Four distinct classes at standard sample points.} The representatives $\{H(0),H(0.4),H(\pi/3),H(2\pi/3)\}$ are pairwise CHM-inequivalent (minimum residual $\gtrsim0.49$; \repo{results/verify\_nwit1\_referee\_checks.txt}).
    Combined with item~(1), the seven values $\{0,0.4,\pi/3,2\pi/3,\pi,4\pi/3,5\pi/3\}$ cover exactly these four CHM classes.
    \item \emph{Not a parametrization gauge.} Distinct $\lambda$ modulo $2\pi$ never yield identical matrices (Lemma~\ref{lem:L4-matrix}); $\lambda$ is a genuine moduli coordinate at the matrix level, even though CHM class is $\pi$-periodic.
\end{enumerate}
In particular, it is \textbf{not} claimed that every pair $\lambda\not\equiv\lambda'\pmod{2\pi}$ gives CHM-inequivalent matrices---that would contradict item~(1).
A complete invariant-theoretic classification of all $\lambda\in[0,2\pi)$ modulo diagonal-phase/permutation CHM moves remains open; items~(1)--(2) record the numerical class structure used for the seven-point endpoint tests below.
\end{lemma}

\begin{proof}
Item~(3) is Lemma~\ref{lem:L4-matrix}.
Items~(1)--(2) are the CHM-equivalence computations recorded in the cited verification logs.
\end{proof}

\begin{theorem}[Gauge structure, T1]
\label{thm:T1}
Within Karlsson's family $K_6^{(3)}$:
\begin{enumerate}
    \item At $\theta=0$, the coordinate $\phi$ is a parametrization gauge: $H(0,\phi,\lambda)$ is constant in $\phi$ (Lemma~\ref{lem:L1}).
    \item The family is $2\pi$-periodic in $\lambda$ (Lemma~\ref{lem:L2}).
    \item At the Dita slice $(\theta_D,\pi/4)$, the block $A$ is the fixed unitary in Lemma~\ref{lem:L3}.
    \item On the Dita slice, $\lambda$ is not a matrix-level gauge (Lemma~\ref{lem:L4-matrix}), while CHM class is $\pi$-periodic with four distinct classes among $\{0,0.4,\pi/3,2\pi/3,\pi,4\pi/3,5\pi/3\}$ (Lemma~\ref{lem:L4}).
\end{enumerate}
\end{theorem}

\begin{proof}
Immediate from Lemmas~\ref{lem:L1}--\ref{lem:L4}.
\end{proof}
